\section{RELATED WORK}
\label{sec:related_work}

The objective of this section is to survey the pioneering research and modern works on the topic with FL in NLP, MTL, and SLM. We map where these fields intersect to allow efficient distributed learning, and the research gaps we propose a real-time framework, specifically.

\subsection{FL for NLP}

Examples of FL in NLP/NLU applications include text classification and personalized language services where raw text cannot be aggregated \cite{mcmahan2017communication, li2020federated_prox}. Fundamental optimization algorithms-such as FedAvg \cite{mcmahan2017communication} which defined the client-update-server-aggregation loop-had done so, and FedProx \cite{li2020federated_prox} performed regularization to enable enhanced behavior to counter client drift. Further improvements, such as SCAFFOLD \cite{karimireddy2020scaffold}, mitigate gradient heterogeneity to guarantee convergence stability. Federated MTL methods like MOCHA illustrate that inter-task relatedness can be strengthened under heterogeneous settings through enhanced statistical efficiency and stability to system issues \cite{smith2017federated}. And newer approaches such as FedBone \cite{chen2024fedbone}, MIRA \cite{elbakary2025mira} and M-Fed \cite{zhou2025multi} take this the step further by separating between shared and task-dependent elements and even specialized encoder-decoder architectures and adapting these into large scale heterogeneous MTL environments. However, several FL-for-NLP studies are still based on very little task diversity or the individual personalization of user at single task.

\subsection{MTL in NLP}

MTL for NLP usually uses one of the two paradigms: hard parameter sharing or soft parameter sharing \cite{smith2017federated, yu2020gradient}. Unlike many of the popular methods such as adversarial sharing which assumes a task head is separate from each model instance, hard sharing utilizes a common encoder for every model instance, which is useful to save the model from overfitting and cost of training. Unlike hard sharing, soft sharing retains task-specific factors, but promotes similarity in some way through regularization and/or through low interaction. Also, parameter-efficient fine-tuning (PEFT), with adapters \cite{houlsby2019parameter} and modular shared frameworks such as AdapterHub \cite{pfeiffer2020adapterhub}, has allowed rapid multi-task adaptation with less overhead. MTL enhances representation learning by transferring useful linguistic data from a task to another. But this also introduces known risks: task imbalance (over-representation of dominant tasks with lower-level tasks) and negative transfer (one task injuring others). These difficulties are especially high in the federated training case, since the client participants involved in the training process are partially responsible, and the balance of distribution of the data over the rounds is not uniform.

\subsection{SLMs for Efficient NLP}

SLMs including DistilBERT \cite{sanh2019distilbert}, MiniLM \cite{wang2020minilm}, and ALBERT \cite{lan2019albert} demonstrate that strict compression and distillation can maintain fair performance at a cost of parameter number in the balance with inference. The common reduction methods are knowledge distillation, parameter sharing, low-rank factorization, and architectural simplification. These resources make those SLMs particularly well suited to edge and distributed learning: transmission overhead is minimized, on-device learning cost on-device is significantly reduced and federated iteration cycles are speeded up. Recent advancements such as EdgeDistilBERT \cite{gholamiangonabadi2025federated} further improve on the state-of-the-art lightweight multi-task FL on edge devices. In multi-task federated systems, SLMs also allow scaling to larger populations of clients and tasks on resource budgets for the system.

\subsection{FL in the presence of heterogeneity in the data}

The non-IID data remains one of the worst bottleneck in the FL process. This can lead to biased updates at local level and cause global optimization failures at global level due to client-dependent language patterns, domain shifts, and label bias. In heterogeneous NLU deployments, existing differences of client tasks (cross-client task diversity) leads to partial representations of client to serve multiple objectives, exacerbating this issue. To mitigate these heterogeneity effects, strategies such as FedBN \cite{li2021fedbn} and low-rank parameterizations like FedPara \cite{nam2021fedpara} have been proposed. As in past work, aggregation without considering the structure of a task will reduce velocity and final generalization, and mixing heterogeneity (to aggregate) will degrade speed. Due to the fact convergence velocity may hinder good generalization and final speed depending on task structure, there have been some proposed techniques, like personalized layers, proximal regularization, gradient conflict reducing and encoder-decoder decomposition which bridged this gap \cite{li2020federated_prox, yu2020gradient, zhou2025multi}.

\subsection{Research Gap}

Recent systematic reviews \cite{khan2024federated} have highlighted the importance of practical NLU systems, despite a strong foundation in existing literature.

\begin{itemize}
    \item No real-time FL frameworks: in the majority of cases, the models are tuned for periodic offline cycles, instead of dynamic adaptation.
    \item Fractional inclusion of FL + MTL + SLM: the majority of contributions consider efficiency, heterogeneity, and multi-task transfer in isolation.
    \item Insufficient heterogeneity analysis in NLU FL: the present tests not much considers how task variance, data skew, and communication limitations interact.
\end{itemize}

This deficiency paves the way for a federated multi-task approach unified in real time and that enables efficient modeling of language but also includes appropriate optimization of distributed NLU with heterogeneity-awareness.